% TOP_PROBLEM: {"id": 3372, "title": "Are all Fermat Numbers square-free?", "category_id": 1, "category": {"id": 1, "name": "number_theory", "display_name": "Number Theory", "description": "Properties of integers, prime numbers, Diophantine equations.", "slug": "number-theory", "order_index": 1, "created_at": "2026-07-31T15:26:25.670Z"}, "status": "open", "problem_number": "OPG-55810", "set_id": 12}
% FIRST_POSTED: 2026-10-01
% ENTRY_KIND: statement_only
% UnsolvedMath ID 3372; source catalogue code OPG-55810
% !TeX program = lualatex
% Definition and sources only; this file is not an LLM solution attempt.
\documentclass[11pt,a4paper]{article}
\usepackage[margin=25mm]{geometry}
\usepackage{fontspec}
\setmainfont{lmroman10-regular.otf}[BoldFont=lmroman10-bold.otf,
 ItalicFont=lmroman10-italic.otf,BoldItalicFont=lmroman10-bolditalic.otf]
\usepackage{amsmath,amssymb,mathtools}
\usepackage{unicode-math}
\setmathfont{latinmodern-math.otf}
\AtBeginDocument{\let\setminus\smallsetminus}
\usepackage{xurl,hyperref}
\hypersetup{colorlinks=true,urlcolor=blue}
\setcounter{secnumdepth}{0}
\setlength{\parindent}{0pt}
\setlength{\parskip}{5pt}
\setlength{\emergencystretch}{3em}
\begin{document}
\section*{Are all Fermat Numbers square-free?}\label{3372}
{\small UnsolvedMath ID 3372 · OPG-55810 · Number Theory · OpenGarden}
\subsection{Definitions and mathematical statement}
For every integer $n\ge0$, the $n$th Fermat number is
$F_n=2^{2^n}+1$. A positive integer $N$ is square-free if no square
of a prime divides it, equivalently if every prime in its prime
factorization occurs with exponent one. The question is whether
\[
 \forall n\ge0\ \forall\text{ primes }p,\qquad p^2\nmid F_n.
\]
Here $n$ need not be prime and there is no upper bound on it. The
numbers $F_n$ are odd, so only odd primes can divide them. One may
equivalently ask whether $v_p(F_n)\le1$ for all $p,n$, where
$v_p(N)$ is the largest integer $a\ge0$ such that $p^a\mid N$.

Square-freeness permits a Fermat number to have many distinct prime
factors; it does not require the number itself to be prime. The
indices $n=0,1$ give $3$ and $5$ and are included. Exhibiting many
distinct factors, or showing that different Fermat numbers are
coprime, does not exclude repeated factors within one $F_n$.
A counterexample consists of a specific index $n$ and a prime $p$
with $2^{2^n}\equiv-1\pmod{p^2}$, certified by exact arithmetic.
\subsection{Short English statement}
Does every Fermat number have no repeated prime factor?
\subsection{Sources}
\begin{itemize}
\item[S1] ULAM.AI and the UnsolvedMath Contributors, supplied problems.json, record 3372 (OPG-55810), OpenGarden collection. Catalogue identity and source transcription; CC BY 4.0.
\url{https://huggingface.co/datasets/ulamai/UnsolvedMath}.
\item[S2] Open Problem Garden, Are all Fermat Numbers square-free?. Original problem and discussion; the mathematical formulation below is expanded from the supplied source record.
\url{http://www.openproblemgarden.org/op/are_all_fermat_numbers_square_free}.
\end{itemize}
\end{document}
